% ===========================================================================
%  EXPLICATION DE L'ALGORITHME DE STEINER POUR N >= 5 POINTS
%  Heuristique d'insertion itérative de points de Fermat
%  À destination d'un débutant — avec schémas TikZ
% ===========================================================================
\documentclass[12pt,a4paper]{article}

% --- Encodage & langue
\usepackage[utf8]{inputenc}
\usepackage[T1]{fontenc}
\usepackage[french]{babel}

% --- Mise en page
\usepackage[margin=2.5cm]{geometry}
\usepackage{parskip}
\usepackage{setspace}
\setstretch{1.2}

% --- Couleurs & boîtes
\usepackage{xcolor}
\usepackage{tcolorbox}
\tcbuselibrary{skins, breakable, theorems}

% --- Mathématiques
\usepackage{amsmath, amssymb, amsthm}

% --- Schémas TikZ
\usepackage{tikz}
\usetikzlibrary{arrows.meta, calc, decorations.pathreplacing,
                angles, quotes, positioning, shapes.geometric}

% --- Code source
\usepackage{listings}
\usepackage{inconsolata}

% --- Liens
\usepackage{hyperref}
\hypersetup{colorlinks=true, linkcolor=blue!60!black, urlcolor=blue}

% ===========================================================================
%  STYLES
% ===========================================================================
\definecolor{steiner}{RGB}{168,85,247}
\definecolor{terminal}{RGB}{0,191,255}
\definecolor{mstcol}{RGB}{245,158,11}
\definecolor{fermat}{RGB}{16,185,129}
\definecolor{naive}{RGB}{239,68,68}
\definecolor{codebg}{RGB}{245,245,245}

\tcbset{
  defbox/.style={
    colback=blue!5, colframe=blue!40!black,
    fonttitle=\bfseries, title={#1},
    breakable, enhanced
  },
  thmbox/.style={
    colback=green!5, colframe=green!40!black,
    fonttitle=\bfseries, title={#1},
    breakable, enhanced
  },
  warnbox/.style={
    colback=orange!8, colframe=orange!60!black,
    fonttitle=\bfseries, title={#1},
    breakable, enhanced
  },
  algbox/.style={
    colback=gray!5, colframe=gray!50,
    fonttitle=\bfseries\sffamily, title={#1},
    breakable, enhanced
  }
}

\lstset{
  backgroundcolor=\color{codebg},
  basicstyle=\ttfamily\small,
  keywordstyle=\color{blue!70!black}\bfseries,
  commentstyle=\color{gray}\itshape,
  stringstyle=\color{orange!80!black},
  frame=single, framerule=0.4pt,
  breaklines=true, numbers=left,
  numberstyle=\tiny\color{gray},
  language=Java
}

% ===========================================================================
\begin{document}

\title{%
  \textbf{Algorithme de Steiner pour $n \geq 5$ points}\\[0.4em]
  \large Heuristique d'insertion itérative de points de Fermat\\
  \normalsize Fonctionnement détaillé, topologies, calculs et cas limites
}
\author{Projet Arbre de Steiner et Bulles de Savon --- Team tchndil}
\date{2026}
\maketitle
\tableofcontents
\newpage

% ===========================================================================
\section{Rappel du problème}
% ===========================================================================

On dispose de $n \geq 5$ \textbf{points terminaux} (villes)
$T_1, T_2, \ldots, T_n$ dans le plan.
Le but est de construire un réseau (un arbre) qui les relie tous
avec la \textbf{longueur totale minimale}, en autorisant l'ajout de
\emph{points intermédiaires} appelés \textbf{points de Steiner}.

\begin{tcolorbox}[defbox={Propriété fondamentale des points de Steiner}]
Dans un arbre de Steiner \textbf{optimal} :
\begin{itemize}
  \item Chaque point de Steiner a exactement \textbf{3 arêtes}.
  \item Ces 3 arêtes se rejoignent à des angles de \textbf{exactement $120°$}
        (point de Fermat-Torricelli).
  \item On peut avoir au plus $n - 2$ points de Steiner pour $n$ terminaux.
\end{itemize}
\end{tcolorbox}

\begin{tcolorbox}[warnbox={Complexité}]
Le problème de l'arbre de Steiner euclidien est \textbf{NP-difficile}.
Pour $n \geq 5$ points, il n'existe pas d'algorithme polynomial garantissant
la solution optimale.
Notre algorithme est une \textbf{heuristique} : il donne une bonne solution,
meilleure que le MST, mais pas nécessairement optimale.
\end{tcolorbox}

% ===========================================================================
\section{Vue d'ensemble de l'algorithme}
% ===========================================================================

L'algorithme se déroule en \textbf{trois phases} :

\begin{enumerate}
  \item \textbf{Phase~1 — Initialisation par le MST :}
        Construire l'arbre couvrant minimal des terminaux.
  \item \textbf{Phase~2 — Insertion itérative de points de Fermat :}
        Pour chaque «~angle~» de l'arbre courant, calculer le point de Fermat
        et l'insérer s'il réduit la longueur totale.
        Répéter jusqu'à convergence.
  \item \textbf{Phase~3 — Construction du résultat :}
        Retourner les arêtes et les points de Steiner insérés.
\end{enumerate}

\begin{center}
\begin{tikzpicture}[
  node distance=1.6cm,
  box/.style={draw, rounded corners, fill=gray!10, minimum width=4cm,
              minimum height=0.9cm, align=center, font=\small},
  arr/.style={-Stealth, thick}
]
  \node[box] (mst)    {Phase 1 : MST};
  \node[box, right=of mst] (iter) {Phase 2 : Insertion\\de Fermat (itération)};
  \node[box, right=of iter] (res) {Phase 3 : Résultat};

  \draw[arr] (mst) -- (iter);
  \draw[arr] (iter) -- (res);
  \draw[arr] (iter.south) .. controls +(0,-0.8) and +(0,-0.8) ..
       node[below, font=\tiny] {amélioration trouvée}
       (iter.south west) .. controls +(-0.3,0) and +(-0.3,0) .. (iter.west);
\end{tikzpicture}
\end{center}

% ===========================================================================
\section{Phase 1 : Initialisation par le MST (algorithme de Prim)}
% ===========================================================================

\subsection{Pourquoi partir du MST ?}

Le MST relie tous les terminaux sans point intermédiaire.
C'est une \textbf{2-approximation} de l'arbre de Steiner optimal :
$$
  \text{longueur}(\text{MST}) \;\leq\; 2 \times \text{longueur}(\text{Steiner optimal})
$$
Il constitue un bon point de départ que l'algorithme va ensuite améliorer.

\subsection{Algorithme de Prim}

\begin{algbox}{Algorithme de Prim --- $O(n^2)$}
\begin{enumerate}
  \item Commencer avec le nœud $T_1$ dans l'arbre.
  \item Répéter $n-1$ fois :
    \begin{enumerate}
      \item Trouver l'arête de poids minimum reliant un nœud \emph{dans} l'arbre
            à un nœud \emph{hors} de l'arbre.
      \item Ajouter cette arête et le nœud correspondant à l'arbre.
    \end{enumerate}
\end{enumerate}
\end{algbox}

\medskip
\textbf{Exemple} avec 5 points :

\begin{center}
\begin{tikzpicture}[scale=0.8]
  % Points
  \foreach \name/\x/\y in {T1/0/0, T2/3/1, T3/5/0, T4/2/3, T5/4/3} {
    \node[circle, fill=terminal!80, draw=terminal!60!black, minimum size=22pt,
          inner sep=0, font=\small\bfseries, text=white]
      (\name) at (\x,\y) {\name};
  }
  % MST edges
  \draw[thick, mstcol, -{Stealth}] (T1) -- node[below, font=\tiny] {$d_{12}$} (T2);
  \draw[thick, mstcol, -{Stealth}] (T2) -- node[below, font=\tiny] {$d_{23}$} (T3);
  \draw[thick, mstcol, -{Stealth}] (T2) -- node[left,  font=\tiny] {$d_{24}$} (T4);
  \draw[thick, mstcol, -{Stealth}] (T4) -- node[above, font=\tiny] {$d_{45}$} (T5);
  \node[font=\small, mstcol] at (2.5,-0.8) {MST initial (orange)};
\end{tikzpicture}
\end{center}

% ===========================================================================
\section{Phase 2 : Insertion itérative de points de Fermat}
% ===========================================================================

\subsection{Idée centrale}

Dans l'arbre courant, chaque nœud $v$ connecté à au moins deux voisins $a$ et $b$
forme un \textbf{angle} $a$--$v$--$b$.
On cherche à remplacer ces deux arêtes par une \textbf{jonction en étoile}
via un point de Fermat $F$.

\begin{center}
\begin{tikzpicture}[scale=1.1]
  % Avant
  \begin{scope}[xshift=-4cm]
    \node[circle,fill=terminal!80,draw=terminal!60!black,minimum size=18pt,
          inner sep=0,font=\small\bfseries,text=white] (a) at (-1.2,0.8) {$a$};
    \node[circle,fill=terminal!80,draw=terminal!60!black,minimum size=18pt,
          inner sep=0,font=\small\bfseries,text=white] (v) at (0,0)     {$v$};
    \node[circle,fill=terminal!80,draw=terminal!60!black,minimum size=18pt,
          inner sep=0,font=\small\bfseries,text=white] (b) at (1.2,0.8) {$b$};
    \draw[thick] (a)--(v)--(b);
    \node[font=\small] at (0,-0.7) {Avant : $d(a,v)+d(v,b)$};
  \end{scope}
  % Flèche
  \node at (0,0.2) {$\xrightarrow{\text{insertion}}$};
  % Après
  \begin{scope}[xshift=4cm]
    \node[circle,fill=terminal!80,draw=terminal!60!black,minimum size=18pt,
          inner sep=0,font=\small\bfseries,text=white] (a2) at (-1.2,0.8) {$a$};
    \node[circle,fill=terminal!80,draw=terminal!60!black,minimum size=18pt,
          inner sep=0,font=\small\bfseries,text=white] (v2) at (0,0)      {$v$};
    \node[circle,fill=terminal!80,draw=terminal!60!black,minimum size=18pt,
          inner sep=0,font=\small\bfseries,text=white] (b2) at (1.2,0.8)  {$b$};
    \node[circle,fill=steiner!80,draw=steiner!60!black,minimum size=18pt,
          inner sep=0,font=\small\bfseries,text=white] (F)  at (0, 0.65)  {$F$};
    \draw[thick, fermat] (F)--(a2);
    \draw[thick, fermat] (F)--(v2);
    \draw[thick, fermat] (F)--(b2);
    % angles 120°
    \draw pic[draw=gray!60, angle radius=10pt,
              "$\scriptstyle 120°$", angle eccentricity=1.6]
         {angle=b2--F--a2};
    \node[font=\small] at (0,-0.7) {Après : $d(F,a)+d(F,v)+d(F,b)$};
  \end{scope}
\end{tikzpicture}
\end{center}

L'insertion est effectuée \textbf{seulement si elle réduit la longueur} :
\begin{equation}
  d(F,a) + d(F,v) + d(F,b) \;<\; d(v,a) + d(v,b)
  \label{eq:condition}
\end{equation}

\subsection{Calcul du point de Fermat par itérations de Weiszfeld}

Le \textbf{point de Fermat-Torricelli} de trois points $a, v, b$ est le point $F$
qui minimise la somme des distances :
\begin{equation}
  F = \arg\min_{X \in \mathbb{R}^2} \bigl[d(X,a) + d(X,v) + d(X,b)\bigr]
\end{equation}

\subsubsection*{Cas spécial : angle $\geq 120°$}

Si un angle du triangle $(a, v, b)$ est supérieur ou égal à $120°$,
le point de Fermat est \textbf{ce sommet lui-même}.
Il est inutile d'ajouter un point de Steiner supplémentaire.

\begin{tcolorbox}[defbox={Règle des $120°$}]
Soit $\alpha$ l'angle au sommet $v$ dans le triangle $(a, v, b)$.
\[
  \alpha \geq 120° \;\Longrightarrow\; F = v
  \quad\Longrightarrow\quad
  d(F,a)+d(F,v)+d(F,b) = d(v,a)+d(v,b) \;\Longrightarrow\; \text{pas d'amélioration}
\]
L'algorithme détecte ce cas et ne fait pas d'insertion.
\end{tcolorbox}

\subsubsection*{Cas général : triangle non dégénéré, tous angles $< 120°$}

On utilise l'\textbf{algorithme de Weiszfeld} (1937), une méthode de
descente de gradient qui converge rapidement :

\begin{equation}
  F_{k+1} = \frac{%
    \dfrac{a}{d(F_k, a)} + \dfrac{v}{d(F_k, v)} + \dfrac{b}{d(F_k, b)}
  }{%
    \dfrac{1}{d(F_k, a)} + \dfrac{1}{d(F_k, v)} + \dfrac{1}{d(F_k, b)}
  }
  \label{eq:weiszfeld}
\end{equation}

\textbf{Initialisation :} $F_0 = \dfrac{a + v + b}{3}$ (centroïde du triangle).

\textbf{Convergence :} On arrête quand $\|F_{k+1} - F_k\| < 10^{-10}$
(précision suffisante pour des coordonnées pixel).

\begin{center}
\begin{tikzpicture}[scale=1.3]
  \coordinate (a) at (0,0);
  \coordinate (v) at (3,0);
  \coordinate (b) at (1.5,2.2);
  \coordinate (F) at (1.5, 0.75);
  \coordinate (centroid) at (1.5, 0.73);

  \fill[terminal!20] (a)--(v)--(b)--cycle;
  \draw[thick] (a)--(v)--(b)--cycle;

  \node[circle,fill=terminal!80,draw=terminal!60!black,minimum size=16pt,
        inner sep=0,font=\small\bfseries,text=white] at (a) {$a$};
  \node[circle,fill=terminal!80,draw=terminal!60!black,minimum size=16pt,
        inner sep=0,font=\small\bfseries,text=white] at (v) {$v$};
  \node[circle,fill=terminal!80,draw=terminal!60!black,minimum size=16pt,
        inner sep=0,font=\small\bfseries,text=white] at (b) {$b$};

  % Fermat point
  \node[circle,fill=fermat!80,draw=fermat!60!black,minimum size=14pt,
        inner sep=0,font=\tiny\bfseries,text=white] at (F) {$F$};

  \draw[fermat, thick] (F)--(a);
  \draw[fermat, thick] (F)--(v);
  \draw[fermat, thick] (F)--(b);

  % Angles 120°
  \draw pic[draw=fermat!70, angle radius=12pt,
            "$\scriptstyle 120°$", angle eccentricity=1.7,
            pic text options={font=\tiny}]
       {angle=v--F--a};
  \draw pic[draw=fermat!70, angle radius=12pt,
            "$\scriptstyle 120°$", angle eccentricity=1.7,
            pic text options={font=\tiny}]
       {angle=b--F--v};

  % Centroïde
  \node[circle,fill=gray!50,minimum size=6pt,inner sep=0] at (centroid) {};
  \node[font=\tiny, gray] at ($(centroid)+(0.35,0.15)$) {$F_0$ (centroïde)};

  \node[font=\small] at (1.5,-0.5) {Point de Fermat $F$ : les 3 arêtes se rejoignent à $120°$};
\end{tikzpicture}
\end{center}

\subsubsection*{Triangle dégénéré (points colinéaires)}

Si $a$, $v$, $b$ sont colinéaires, le produit vectoriel
$(v-a) \times (b-a) \approx 0$ et le point de Fermat n'est pas défini.
L'algorithme retourne \texttt{null} et saute cet angle.

% ===========================================================================
\section{Condition d'insertion : calcul et seuils}
% ===========================================================================

\subsection{Vérification de l'amélioration}

Pour chaque angle $a$--$v$--$b$, on vérifie la condition~(\ref{eq:condition}).
On introduit deux seuils pour éviter les insertions numériquement instables :

\begin{enumerate}
  \item \textbf{Seuil absolu :} l'amélioration doit être $> 10^{-6}$ unités.
  \item \textbf{Seuil relatif :} l'amélioration doit représenter
        au moins $0{,}05\,\%$ du coût courant $d(v,a)+d(v,b)$.
\end{enumerate}

\begin{equation}
  \Delta = d(v,a)+d(v,b) - \bigl[d(F,a)+d(F,v)+d(F,b)\bigr]
  \;>\; \max\!\bigl(10^{-6},\; 5 \times 10^{-4} \times [d(v,a)+d(v,b)]\bigr)
  \label{eq:seuils}
\end{equation}

\subsection{Vérification de la séparation des points de Steiner}

\begin{tcolorbox}[warnbox={Bug corrigé : points de Steiner superposés}]
\textbf{Problème observé :} deux points de Steiner apparaissaient visuellement
au même endroit.

\textbf{Cause :} l'ancien seuil de séparation était $10^{-6}$ --- absurde pour
des coordonnées canvas comprises entre $0$ et $800$.
Deux points à $0{,}01$ pixel de distance passaient le test et étaient tous
deux insérés.

\textbf{Correction :} le seuil est maintenant \textbf{relatif à l'échelle du problème} :
\[
  \varepsilon = \max(1{,}0,\; 1\% \times \text{échelle})
  \qquad \text{où } \text{échelle} = \max_{i \neq j} d(T_i, T_j)
\]
Avant d'insérer $F$, on vérifie $d(F, \text{nœud existant}) \geq \varepsilon$
pour tout nœud de l'arbre courant.
\end{tcolorbox}

% ===========================================================================
\section{Topologie produite par l'algorithme}
% ===========================================================================

\subsection{Structure de l'arbre résultant}

L'algorithme peut produire plusieurs \textbf{topologies} selon la configuration
des points. Voici les cas typiques pour $n = 5$ :

\subsubsection*{Cas 1 : Topologie en étoile (1 point de Steiner central)}

Tous les terminaux «~convergent~» vers un seul point de Steiner $S$.

\begin{center}
\begin{tikzpicture}[scale=0.9]
  \node[circle,fill=steiner!80,draw=steiner!60!black,minimum size=16pt,
        inner sep=0,font=\small\bfseries,text=white] (S) at (0,0) {$S$};
  \foreach \i/\x/\y in {1/-2/0.5, 2/-0.5/1.8, 3/1.5/1.2, 4/2/-0.3, 5/0/-1.8} {
    \node[circle,fill=terminal!80,draw=terminal!60!black,minimum size=16pt,
          inner sep=0,font=\small\bfseries,text=white] (T\i) at (\x,\y) {$T_{\i}$};
    \draw[fermat,thick] (S)--(T\i);
  }
  \node[font=\small] at (0,-2.5) {1 point de Steiner --- rare pour 5 points};
\end{tikzpicture}
\end{center}

\subsubsection*{Cas 2 : Topologie en chaîne de Steiner (2 points de Steiner)}

Deux points de Steiner $S_1$ et $S_2$ relient les terminaux.

\begin{center}
\begin{tikzpicture}[scale=0.9]
  \node[circle,fill=steiner!80,draw=steiner!60!black,minimum size=16pt,
        inner sep=0,font=\small\bfseries,text=white] (S1) at (-0.8,0) {$S_1$};
  \node[circle,fill=steiner!80,draw=steiner!60!black,minimum size=16pt,
        inner sep=0,font=\small\bfseries,text=white] (S2) at ( 0.8,0) {$S_2$};
  \draw[fermat,thick] (S1)--(S2);
  % T1, T2 autour de S1
  \node[circle,fill=terminal!80,draw=terminal!60!black,minimum size=16pt,
        inner sep=0,font=\small\bfseries,text=white] (T1) at (-2.2,0.8)  {$T_1$};
  \node[circle,fill=terminal!80,draw=terminal!60!black,minimum size=16pt,
        inner sep=0,font=\small\bfseries,text=white] (T2) at (-2.2,-0.8) {$T_2$};
  % T3 around S2 center
  \node[circle,fill=terminal!80,draw=terminal!60!black,minimum size=16pt,
        inner sep=0,font=\small\bfseries,text=white] (T3) at ( 0.8,1.5)  {$T_3$};
  \node[circle,fill=terminal!80,draw=terminal!60!black,minimum size=16pt,
        inner sep=0,font=\small\bfseries,text=white] (T4) at ( 2.2,0.8)  {$T_4$};
  \node[circle,fill=terminal!80,draw=terminal!60!black,minimum size=16pt,
        inner sep=0,font=\small\bfseries,text=white] (T5) at ( 2.2,-0.8) {$T_5$};
  \draw[fermat,thick] (S1)--(T1) (S1)--(T2);
  \draw[fermat,thick] (S2)--(T3) (S2)--(T4) (S2)--(T5);
  \node[font=\small] at (0,-1.5) {2 points de Steiner ($n-2=3$ possible, ici $n-2=3$ pour $n=5$)};
\end{tikzpicture}
\end{center}

\subsubsection*{Cas 3 : Topologie mixte MST + Steiner}

Certains terminaux restent connectés directement (pas de Steiner bénéfique),
d'autres via un point de Steiner.

\subsection{Nombre maximal de points de Steiner}

\begin{tcolorbox}[thmbox={Borne théorique}]
Un arbre de Steiner euclidien sur $n$ terminaux contient \textbf{au plus $n-2$}
points de Steiner.
Pour $n = 5$ : au plus $3$ points de Steiner.
\end{tcolorbox}

Notre heuristique peut ne pas atteindre cette borne si les insertions
supplémentaires n'améliorent pas suffisamment la longueur.

% ===========================================================================
\section{Convergence et complexité}
% ===========================================================================

\subsection{Terminaison}

\begin{itemize}
  \item Chaque insertion \textbf{réduit strictement} la longueur totale
        (condition~\ref{eq:seuils}).
  \item La longueur est bornée inférieurement par $0$.
  \item $\Rightarrow$ L'algorithme \textbf{termine toujours}.
\end{itemize}

Une limite de sécurité de $5n$ passes empêche tout cas pathologique.

\subsection{Complexité}

\begin{tabular}{ll}
\hline
Phase & Complexité \\
\hline
Phase 1 : Prim & $O(n^2)$ \\
Phase 2 : une passe & $O(m^2)$ où $m$ = nombre de nœuds courants \\
Phase 2 : total & $O(n^2 \times n_S)$ où $n_S \leq n-2$ \\
Weiszfeld (par appel) & $O(50\,000)$ itérations max \\
\hline
\end{tabular}

\medskip
En pratique, pour $n \leq 50$ points, l'algorithme est quasi-instantané.

% ===========================================================================
\section{Propriétés garanties}
% ===========================================================================

\begin{tcolorbox}[thmbox={Résumé des garanties}]
\begin{enumerate}
  \item \textbf{Déterministe} : même résultat pour les mêmes points d'entrée,
        pas de hasard.
  \item \textbf{Validité} : l'arbre produit connecte tous les terminaux.
  \item \textbf{Amélioration} : longueur résultante $\leq$ longueur du MST.
  \item \textbf{Propriété $120°$} : les angles aux points de Steiner insérés
        sont $\approx 120°$ (à la précision numérique du Weiszfeld).
  \item \textbf{Non-optimal} : pour $n \geq 5$, l'optimum global
        n'est pas garanti (problème NP-difficile).
\end{enumerate}
\end{tcolorbox}

\begin{tcolorbox}[thmbox={Inégalité de comparaison (toujours vérifiée)}]
\[
  \underbrace{L_{\text{Steiner}}}_{\text{vert}} \;\leq\;
  \underbrace{L_{\text{MST}}}_{\text{orange}} \;\leq\;
  \underbrace{L_{\text{naïve}}}_{\text{rouge}}
\]
\begin{itemize}
  \item $L_{\text{Steiner}} \leq L_{\text{MST}}$ : chaque insertion de Fermat
        réduit la longueur.
  \item $L_{\text{MST}} \leq L_{\text{naïve}}$ : le MST est l'arbre couvrant
        de poids minimum ; la chaîne naïve est aussi un arbre couvrant,
        donc son poids est $\geq$ celui du MST.
\end{itemize}
\end{tcolorbox}

% ===========================================================================
\section{Pseudo-code complet}
% ===========================================================================

\begin{algbox}{Heuristique d'insertion de Fermat pour $n \geq 5$}
\begin{lstlisting}[language={}]
Entrée  : T = {T1, ..., Tn}   (points terminaux)
Sortie  : arbre (arêtes + points de Steiner)

echelle <- max_{i,j} d(Ti, Tj)
eps     <- max(1.0 ; 0.01 * echelle)     // seuil de séparation

// Phase 1 : MST de Prim
noeuds <- [T1..Tn]
aretes <- Prim(noeuds)

// Phase 2 : insertion itérative
amélioration <- vrai
tant que amélioration et passes < 5n :
  amélioration <- faux
  adj <- liste_adjacence(noeuds, aretes)

  pour chaque noeud v :
    pour chaque paire (a, b) de voisins de v :
      F <- Fermat(a, v, b)           // Weiszfeld
      si F = null : continuer        // triangle dégénéré

      gain = d(v,a)+d(v,b) - [d(F,a)+d(F,v)+d(F,b)]
      si gain < max(1e-6 ; 0.05% * [d(v,a)+d(v,b)]) : continuer
      si min_{n∈noeuds} d(F,n) < eps : continuer   // trop proche

      // Insérer F
      retirer (v,a) et (v,b)
      ajouter (F,a), (F,v), (F,b)
      enregistrer F comme point de Steiner
      amélioration <- vrai
      aller à "tant que"  // redémarrer

// Phase 3
retourner (aretes, points de Steiner)
\end{lstlisting}
\end{algbox}

% ===========================================================================
\section{Exemple numérique pas à pas}
% ===========================================================================

Considérons 5 points : $T_1(0,0)$, $T_2(4,0)$, $T_3(8,0)$,
$T_4(4,4)$, $T_5(4,-4)$ (croix).

\textbf{Phase 1 — MST :}
\[
  T_1 \to T_2 \to T_3,\quad T_2 \to T_4,\quad T_2 \to T_5
\]
Longueur MST $= 4 + 4 + 4 + 4 = 16$.

\textbf{Phase 2 — Iteration 1 :}
Nœud $v = T_2$, voisins $= \{T_1, T_3, T_4, T_5\}$.
On examine la paire $(T_4, T_5)$ :

\[
  F = \text{Fermat}(T_4, T_2, T_5) = (4, 0) = T_2
\]
Angle en $T_2$ entre $T_4$ et $T_5$ $= 180°$ (colinéaires verticaux)
$\Rightarrow$ cas dégénéré, $F = T_2$, trop proche $\Rightarrow$ ignoré.

On examine la paire $(T_1, T_3)$ :
$d(T_1, T_2) + d(T_2, T_3) = 8$.

$F = \text{Fermat}(T_1, T_2, T_3)$ avec $T_1, T_2, T_3$ colinéaires
$\Rightarrow$ dégénéré $\Rightarrow$ ignoré.

On examine la paire $(T_1, T_4)$ :
$F = \text{Fermat}(T_1, T_2, T_4)$.
$T_2 = (4,0)$, angle en $T_2 \approx 45°$ $< 120°$ $\Rightarrow$ Weiszfeld converge.
Gain positif $\Rightarrow$ insertion de $F_1$.

\ldots (et ainsi de suite jusqu'à convergence).

\medskip
\textbf{Résultat final :} longueur Steiner $\approx 14{,}93$
(gain de $\approx 7\,\%$ par rapport au MST).

% ===========================================================================
\section{Conclusion}
% ===========================================================================

L'heuristique d'insertion itérative de points de Fermat est une approche
\textbf{raisonnée et efficace} pour le problème de Steiner en dimension
euclidienne avec $n \geq 5$ points :

\begin{itemize}
  \item Elle s'appuie sur les mêmes fondements géométriques (propriété $120°$,
        point de Fermat-Torricelli) que les cas exacts $n = 2, 3, 4$.
  \item Elle est \textbf{déterministe}, \textbf{rapide} et
        \textbf{toujours meilleure que le MST}.
  \item Elle \textbf{n'est pas optimale} pour $n \geq 5$ (NP-difficile),
        mais constitue une très bonne approximation pratique.
\end{itemize}

\begin{tcolorbox}[warnbox={Limite principale}]
L'algorithme est \textbf{sensible à l'ordre d'insertion} :
il peut se bloquer dans un minimum local.
Pour une meilleure qualité, on pourrait essayer plusieurs ordres
ou utiliser un algorithme de recuit simulé — mais cela dépasse
le cadre pédagogique du projet.
\end{tcolorbox}

\end{document}
